%% Adequacy and decision links, with remaining sketches explicitly identified.

\begin{annexitem}{annex:p8.ipl.source}{entry:p08-ipl}{IPL soundness, strong completeness and finite countermodels}
\begin{proof}
The card uses exactly the common-world consequence and finite-support
propositional LJ derivability of
Definition~\ref{def:p08-ipl-completeness-setting}. Rule-by-rule soundness is
Theorem~\ref{thm:p08-ipl-completeness-soundness}. The prime-extension and
failed-implication constructions and canonical truth lemma prove
Theorem~\ref{thm:p08-ipl-strong-completeness-full} for every premise set in
classical ZFC. For finite inputs,
Theorem~\ref{thm:p08-ipl-finite-rooted-countermodel} supplies the rooted finite
countermodel by the fully proved filtration. These results prove both claims
in the adequacy field, with the same language and frame class.
\end{proof}
\end{annexitem}

\begin{annexitem}{annex:p8.ipl.decision-card}{entry:p08-ipl}{IPL finite-input decision}
\begin{proof}
Theorem~\ref{thm:p08-ipl-finite-input-decision} bounds a countermodel by the
finite subformula profiles of the input. Enumerating the finite rooted
partial orders and persistent valuations within that bound and evaluating
each formula is a terminating finite procedure. Soundness and the finite
countermodel theorem prove the two possible answers correct. This decides
finite-input propositional derivability, not effectively presented arbitrary
infinite premise sets.
\end{proof}
\end{annexitem}

\begin{annexitem}{annex:p8.modal.source}{thm:p08-modal-source-sketch}{Inherited four-system modal completeness sketch}
\sketchnote{Lindenbaum extension, the modal existence lemma, the full truth
lemma, and the separate T, 4, and 5 canonical-frame calculations are compressed
in the inherited sketch.}
The theorem and reader-facing sketch remain together at
\hyperref[thm:p08-source-modal-four-adequacy]{the exact inherited
frame-correspondence result}.  It defines the canonical relation, identifies
the truth-lemma induction, and names the frame-property consequences of T, 4,
and 5.  This item records that standing; the separately authored complete
argument is Annex~\ref{annex:p8.modal.1}.
\end{annexitem}

\begin{annexitem}{annex:p8.ltl.source}{entry:p08-ltl}{Infinite-word LTL soundness and weak completeness}
\begin{proof}
The card is the registered next/non-strict strong-until language.
Theorem~\ref{thm:p08-ltl-soundness} proves soundness of its seven Hilbert
rules. The consistent-type, successor, eventuality and truth-lemma
constructions prove Theorem~\ref{thm:p08-linear-temporal-completeness}.
Underivability yields an infinite word falsifying the formula, while
derivability gives truth at every position of every word. Neither the
strict-future source variant nor arbitrary-premise strong completeness is
needed for this equivalence.
\end{proof}
\end{annexitem}

\begin{annexitem}{annex:p8.ltl.decision-card}{entry:p08-ltl}{LTL satisfiability and combined model-checking complexity}
\begin{proof}
The complete accepting-cycle procedure and
Theorem~\ref{thm:p08-ltl-polynomial-space} give deterministic polynomial-space
upper bounds using the explicit encoded input lengths. The machine-history
construction of Theorem~\ref{thm:p08-ltl-pspace-hardness} proves satisfiability
hardness. The block translation and finite total transition model of
Theorem~\ref{thm:p08-ltl-model-checking-complete} prove the matching combined
model-checking classification. The universal-path check negates the
specification before testing an accepting product path. These are the two
PSPACE-complete problems stated by the card, not a data-complexity claim.
\end{proof}
\end{annexitem}

\begin{annexitem}{annex:p8.s5n.source}{entry:p08-s5n}{Inherited multi-agent S5 adequacy sketch}
\sketchnote{Lindenbaum extension, the modal existence lemma, and the full
truth-lemma calculation are explicitly omitted by the inherited sketch.}
The theorem and reader-facing sketch remain together at
\hyperref[thm:p08-source-s5n-adequacy]{the exact inherited multi-agent S5
result}.  The printed argument checks the componentwise S5 axioms, defines
each canonical relation, and verifies reflexivity, transitivity, and
Euclideanness without adding cross-agent interaction; this annex item is only
its reciprocal route.
\end{annexitem}

\begin{annexitem}{annex:p8.imll.source}{entry:p08-imll}{Inherited IMLL adequacy and cut-elimination sketch}
\sketchnote{The rule cases for soundness, the Lindenbaum quotient details, and
the principal and commuting cut reductions are compressed rather than proved
case by case in the inherited sketch.}
The theorem and reader-facing sketch remain together at
\hyperref[thm:p08-source-imll-adequacy]{the exact inherited IMLL adequacy
result}.  The printed argument identifies the monoid and residual in the
Lindenbaum quotient and names the context-splitting cut reductions; this annex
item only supplies reciprocal navigation.
\end{annexitem}
